\documentclass[10pt, aspectratio=169]{beamer}
\usepackage[utf8]{inputenc}
\usepackage[spanish,es-nodecimaldot]{babel}
\usepackage{amsmath, amssymb}
\usepackage{booktabs}
\usepackage{tikz}
\usetikzlibrary{shapes.misc, positioning, arrows.meta}

% Configuración de colores y estilo visual
\usetheme{Madrid}
\usecolortheme{seahorse}

\definecolor{mainblue}{RGB}{40, 60, 110}
\definecolor{lightbg}{RGB}{240, 243, 250}
\definecolor{blockbg}{RGB}{230, 235, 245}

\setbeamercolor{structure}{fg=mainblue}
\setbeamercolor{title}{fg=mainblue, bg=lightbg}
\setbeamercolor{frametitle}{fg=mainblue, bg=lightbg}
\setbeamercolor{block title}{fg=white, bg=mainblue}
\setbeamercolor{block body}{fg=black, bg=blockbg}

% Título y Metadatos
\title[Laboratorio 3 sEMG reproducible]{Laboratorio 3 - Dashboard interactivo y análisis reproducible de sEMG}
\author[Integrantes]{Cristian Stiven Capera C., Ana Sofía García G., Jeimmy Andrea Gonzáles G., Karol Mariana Gutiérrez G. y Brayan Andrés Ruiz C.}
\date{01 de Octubre de 2026}

\begin{document}

% -----------------------------------------------------------------------------
% Diapositiva 1
% -----------------------------------------------------------------------------
\begin{frame}
\titlepage
\begin{columns}[t]
    \begin{column}{0.48\textwidth}
        \begin{block}{Problema y contexto}
        \fontsize{8.5pt}{10.5pt}\selectfont
        La sEMG observa la activación muscular pero depende de factores anatómicos, electrodos y ruido. Interpretar señales crudas puede ocultar decisiones críticas de procesamiento.
        \end{block}
        
        \vspace{0.1cm}
        \fontsize{8.5pt}{10.5pt}\selectfont
        \textbf{Pregunta:} ¿Cómo transformar sEMG dinámica en métricas interpretables sin ocultar las decisiones metodológicas?
    \end{column}
    
    \begin{column}{0.48\textwidth}
        \fontsize{8.5pt}{10.5pt}\selectfont
        \textbf{Usuario:} profesional de rehabilitación física que necesita explorar patrones cuantitativos de activación muscular.
        
        \vspace{0.2cm}
        \textbf{Límite clínico:} El dashboard \textbf{no diagnostica} fatiga, lesión ni sobreesfuerzo. Las métricas son indicadores exploratorios.
    \end{column}
\end{columns}
\end{frame}

% -----------------------------------------------------------------------------
% Diapositiva 2
% -----------------------------------------------------------------------------
\begin{frame}{BTS TDF, arquitectura y configuración}
\begin{columns}[c]
    \begin{column}{0.35\textwidth}
        \centering
        \begin{tikzpicture}[
            node distance=0.3cm,
            block/.style={draw=mainblue!80, fill=lightbg, rectangle, rounded corners=3pt, inner sep=3pt, text width=3.2cm, align=center, font=\fontsize{7pt}{8pt}\selectfont},
            arrow/.style={-Stealth, thick, mainblue}
        ]
            \node[block] (step1) {Archivos TDF\\ \texttt{data/raw/}};
            \node[block, below=of step1] (step2) {\texttt{io.py} + Contrato};
            \node[block, below=of step2] (step3) {\texttt{data\_validation.py}\\Calidad y muestreo};
            \node[block, below=of step3] (step4) {\texttt{preprocessing\_metrics.py}\\DC + banda + envolvente};
            \node[block, below=of step4] (step5) {\texttt{metrics.py}\\RMS + MDF + activaciones};
            \node[block, below=of step5] (step6) {\texttt{app.py} (Streamlit)\\Dashboard y visuales};
            
            \draw[arrow] (step1) -- (step2);
            \draw[arrow] (step2) -- (step3);
            \draw[arrow] (step3) -- (step4);
            \draw[arrow] (step4) -- (step5);
            \draw[arrow] (step5) -- (step6);
        \end{tikzpicture}
    \end{column}
    
    \begin{column}{0.62\textwidth}
        \fontsize{8.5pt}{10.5pt}\selectfont
        \textbf{Registros utilizados}
        \begin{itemize}
            \item Tres registros BTS TDF de flexiones de brazos (61--65 s).
            \item 8 canales sEMG (bíceps, deltoides, pectoral y tríceps bilateral).
            \item $f_s = 1000$ Hz leída directamente del bloque EMG ($f_N = 500$ Hz).
        \end{itemize}
        
        \vspace{0.1cm}
        \textbf{Configuración central}
        \vspace{0.05cm}
        
        \centering
        \fontsize{7.5pt}{9.5pt}\selectfont
        \begin{tabular}{ll}
            \toprule
            \textbf{Parámetro} & \textbf{Valor} \\
            \midrule
            Filtro pasa-banda & 20--450 Hz, orden 4 \\
            Envolvente pasa-bajas & 5 Hz, orden 4 \\
            Ventana RMS móvil & 200 ms \\
            Ventana MDF (Welch) & 2 s (solapamiento 50\%) \\
            Línea base activación & 1 s ($k=3$, MAD) \\
            Duración mínima activ. & 0.10 s \\
            \bottomrule
        \end{tabular}
    \end{column}
\end{columns}
\end{frame}

% -----------------------------------------------------------------------------
% Diapositiva 3
% -----------------------------------------------------------------------------
\begin{frame}{Preprocesamiento, métricas y segmentación}
\fontsize{8.5pt}{10.5pt}\selectfont
\begin{columns}[t]
    \begin{column}{0.48\textwidth}
        \textbf{Pipeline de señal}
        \begin{itemize}
            \item \textbf{DC:} Centrado restando la media.
            \item \textbf{Filtro:} Butterworth de fase cero (\texttt{sosfiltfilt}) 20--450 Hz.
            \item \textbf{Rectificación y envolvente:} Valor absoluto seguido de LPF a 5 Hz.
            \item \textbf{Notch 60 Hz:} Configurable pero desactivado por defecto para evitar distorsión espectral innecesaria.
        \end{itemize}
        
        \vspace{0.1cm}
        \centering
        $RMS = \sqrt{\frac{1}{N}\sum x_n^2}, \quad iEMG = \frac{\sum |x_n|}{f_s}$
    \end{column}
    
    \begin{column}{0.48\textwidth}
        \textbf{Detección de activaciones}
        \begin{itemize}
            \item Umbral robusto basado en el primer segundo:
            $T = \text{med} + 3(1.4826)\text{MAD}$.
            \item Duración mínima de 100 ms para eliminar fluctuaciones fugaces.
            \item Brecha de fusión de 100 ms para evitar fragmentación.
        \end{itemize}
        
        \vspace{0.1cm}
        \textbf{Vector de características}
        \begin{itemize}
            \item RMS, \%MVC, MDF, iEMG, WL, número de activaciones y duración media.
        \end{itemize}
    \end{column}
\end{columns}
\end{frame}

% -----------------------------------------------------------------------------
% Diapositiva 4
% -----------------------------------------------------------------------------
\begin{frame}{Auditoría de parámetros manuales}
\fontsize{8pt}{10pt}\selectfont
Justificación de los parámetros fijados en código (evita presentarlos como constantes fisiológicas universales).

\vspace{0.2cm}
\centering
\begin{tabular}{llll}
    \toprule
    \textbf{Parámetro} & \textbf{Valor} & \textbf{Tipo} & \textbf{Justificación principal} \\
    \midrule
    Frecuencia ($f_s$) & 1000 Hz & Adquisición & Leída del TDF; Nyquist a 500 Hz. \\
    Pasa-banda & 20--450 Hz & Literatura & Reduce movimiento; 450 Hz deja margen vs Nyquist. \\
    Envolvente LPF & 5 Hz & Fisiología & Sigue la escala temporal de flexiones dinámicas. \\
    RMS móvil & 200 ms & Empírico & Equilibra estabilidad y resolución temporal. \\
    Mediana (MDF) & 2 s & Espectral & $\Delta f \approx 0.5$ Hz; menor variabilidad que 1 s. \\
    \bottomrule
\end{tabular}

\vspace{0.4cm}
\begin{block}{Nota metodológica}
El análisis de sensibilidad interno (ventanas RMS de 100/200/400 ms y MDF de 1/2/4 s) respalda estas elecciones como compromisos de ingeniería adecuados para este dataset.
\end{block}
\end{frame}

% -----------------------------------------------------------------------------
% Diapositiva 5
% -----------------------------------------------------------------------------
\begin{frame}[fragile]{Dashboard, reproducibilidad y límites}
\fontsize{8.5pt}{10.5pt}\selectfont
\begin{columns}[t]
    \begin{column}{0.32\textwidth}
        \textbf{Dashboard Streamlit}
        \begin{enumerate}
            \item Carga de TDF o sintético.
            \item Control de filtros y envolvente.
            \item Tendencias RMS y MDF.
            \item Espectro PSD de Welch.
            \item Calidad y activaciones.
        \end{enumerate}
    \end{column}
    
    \begin{column}{0.35\textwidth}
        \textbf{Reproducción}
        \begin{itemize}
            \item Datos originales en \texttt{data/raw/}.
            \item \texttt{scripts/reproduce.py} exporta CSV derivados y JSON.
            \item Pruebas unitarias en \texttt{pytest}.
        \end{itemize}
        
        \vspace{0.1cm}
        \textbf{Comandos}
        \begin{tiny}
\begin{verbatim}
uv sync --locked
uv run ruff check
uv run pytest -v
uv run python scripts/reproduce.py
uv run streamlit run app.py
\end{verbatim}
        \end{tiny}
    \end{column}
    
    \begin{column}{0.31\textwidth}
        \textbf{Límites}
        \begin{itemize}
            \item Unidad física adoptada como V (por confirmar).
            \item Sin MVC real integrada por defecto.
            \item Filtro offline (\texttt{sosfiltfilt}).
            \item Herramienta de apoyo cuantitativo, no diagnóstico clínico.
        \end{itemize}
    \end{column}
\end{columns}
\end{frame}

\end{document}